\scanpage{395}\mspage{134}

qui est de degré $1$, donc qui est une $\mathrm{Hot}$ équivalence. \deleted{Donc} \inserted{Mais} l'application correspondante

\[
\pi_0\underline{\mathrm{Hom}}(Y,\Omega(X,a))\unsure{-|\longrightarrow}\pi_0\underline{\mathrm{Hom}}(Y',\Omega(X,a))
\]

n'est pas bijective, Ce genre d'argument me semble devoir marcher chaque fois que l'un des groupes d'homotopie \deleted{\unsure{du}} $\pi_j(X,x)$, pour $j\geq2$ (donc l'un des groupes d'homotopie $\pi_i$ de $\Omega(X,x)$, pour $i\geq1$) est \uline{infini}, en prenant pour $Y,Y'$ \deleted{\unsure{deux}} différents représentants \inserted{$Y_n$} de la sphère $S^i$, \deleted{\unsure{suivant}} disons \deleted{de la forme $\unsure{\mathscr L|\mathcal L}_n$} les $Y_n$ déduits des \deleted{\unsure{la qu}} $\mathscr L_n$ \inserted{$(n\geq2)$} en prenant $\mathscr L_n^i$ \inserted{$=\mathscr L_n\times\mathscr L_n\,\ldots\times\mathscr L_n$} (qui est une sorte de $i$-cube combinatoire), et en identifiant son ``bord'' à un seul point. Je présume que le nombre d'éléments \inserted{de $\pi_i(\Omega(X,x))$} qu'on peut obtenir à partir de \deleted{l'application} flèches $Y_n\longrightarrow\Omega(X,x)$ est limité par la taille de $Y_n$, mais on peut \inserted{sûrement} les obtenir chacun, en prenant $n$ assez grand. Si donc $\unsure{\zeta|\xi}\in\pi_i(\Omega)$ échappe à $Y_n$ $(n_0=\unsure{n|0}$, disons), on peut \unsure{un|en} $n'$ assez grand pour que \deleted{\unsure{il}} $\zeta$ n'échappe pas à $Y_{n'}$, et \unsure{on|en} choisit

\[
Y_{n'}\longrightarrow Y_n
\]

% archive sheet number: 394
